\section*{Methods}

\subsection*{Configuration space and peg preservation}
At a point in time a stablecoin is a tuple $s_i=(C_i,O_i,R_i)$: a collateral composition $C_i$ in the $(K-1)$-simplex over $K$ reserve asset classes, an oracle price vector $O_i\in[0.5,1.5]^J$ with consensus value $\bar O_i=\phi_i(O_i)$, and a redemption gate state $R_i\in\{\text{open},\text{gated},\text{halted}\}$. The signed peg deviation is $\delta_i:=\bar O_i-1$; the backing ratio $\rho_i$ is derived from the collateral and the issuer's outstanding supply and is not an independent coordinate. A state is peg-preserving at tolerance $\varepsilon$ if $|\delta_i|\le\varepsilon$ and the redemption arbitrage cycle is executable with round-trip profitability bounded by $\varepsilon$; the default $\varepsilon=50$ basis points follows \citep{LyonsViswanath2023}. The threshold enters only the definition of peg preservation and is not used as a cut-off anywhere in the estimation.

The configuration space is $\mathcal{M}=\{(s_1,\dots,s_N):\text{joint feasibility}\}$, where joint feasibility comprises backing balance, oracle consistency and redemption clearing; the peg-preserving locus $\mathcal{M}_{\mathrm{peg}}\subset\mathcal{M}$ is the subset on which every coin is peg-preserving. We assume both are non-empty on the parameter range considered. Peg preservation is a constraint cutting a locus out of the ambient space, not a restriction on what the ambient space may contain: de-pegging configurations remain points of $\mathcal{M}$, and the chambers and walls stratify that ambient space.

\subsection*{The three assumptions}
Write $\Theta_{\mathrm{peg}}$ for the parameter range, with coordinates comprising the structural parameters (issuance rule, collateral haircut, oracle bandwidth) together with the state functionals $\rho_i(\theta)$ and $\delta_i(\theta)$ through which those parameters act, so that a condition such as $\rho_i=1$ is read as the level set $\{\theta:\rho_i(\theta)=1\}$ and defines a wall. The wall set is $\mathcal{W}=\bigcup_{k=1}^{8}W_k$; a chamber is a connected component of $\Theta_{\mathrm{peg}}\setminus\mathcal{W}$. Four walls are equalities between a continuous parameter and a threshold (backing ratio, oracle deviation, funding basis, interest spread) and are codimension-one hypersurfaces; the other four are discrete transitions (redemption gate, custody, bridge, regulation). All eight are retained in one list because the crossing each induces has the same rank-one form, but only the first group can be drawn as boundaries in a continuous parameter plane.

For a chamber and a parameter $\theta$ in it, let $\mu_\theta$ be the probability law of $\delta$ while the system sits in that chamber. The law is a primitive of the framework: no dynamics is written down, and the only properties used are that $\mu_\theta$ is a probability measure with finite second moments. The co-movement covariance is the centred second moment
\[
C_{ij}(\theta):=\operatorname{Cov}_{\mu_\theta}(\delta_i,\delta_j),\qquad
C(\theta):=\mathbb{E}_{\mu_\theta}\bigl[(\delta-\bar\delta)(\delta-\bar\delta)^\top\bigr],
\]
which is symmetric and positive semidefinite by construction. The centring is not cosmetic: a coin that trades persistently below par has $\bar\delta_i\ne0$, and an uncentred second moment would not be recovered by any correlation estimate computed on its price series. The normalised image is $\Sigma_{ij}=C_{ij}/\sqrt{C_{ii}C_{jj}}$ on the open set where $C_{ii}>0$, with $\Sigma_{ii}:=1$.

\textbf{(H1) Within-chamber rigidity.} For any two parameter values in the same chamber there is a scalar $\lambda>0$ with $C(\theta)=\lambda C(\theta')$. The scale factor allows the overall level of volatility to drift within a regime while the shape of the co-movement structure is fixed.

\textbf{(H2) Rank-one wall crossing.} Across a wall $W$ separating two adjacent chambers,
\[
C^{+}-C^{-}=\sigma_W\sum_{\gamma\in\Gamma_W}\Omega(\gamma)\,v_\gamma v_\gamma^\top ,
\]
with $\Gamma_W$ a finite set of channels active at $W$, $\Omega(\gamma)\in\mathbb{R}$ signed weights, $v_\gamma$ directions supported on the coins channel $\gamma$ affects, and $\sigma_W>0$ a positive scale. The form is that of the Kontsevich--Soibelman wall-crossing formula \citep{KontsevichSoibelman2008}, which takes a product over the charges whose central charges change sign across a wall and, where it applies, yields a signed sum of rank-one terms on the coins those charges carry. We take that form as the primitive. Two features of the original are not transported: $\Omega$ is not claimed to be an integer, since only the product $\sigma_W\Omega$ is ever identified and integrality would carry no empirical content; and no sign restriction is placed on $\Omega$, so the right-hand side is a general symmetric matrix. The crossing law therefore fixes the support of the jump but not its direction.

\textbf{(H3) Block support.} For each wall crossed by an episode with a multi-coin channel, the directions $\{v_\gamma\}$ span the indicator directions of a single group of coins $B_W$, the wall-affected block; coins outside $B_W$ carry zero weight. $B_W$ is fixed by the channel's structural definition and never by the observed correlation outcomes.

\subsection*{Consequences}
Because $\Sigma_{ii}=1$ on both sides of a wall, the induced correlation jump $\Delta_R:=\Sigma^{+}-\Sigma^{-}$ has zero diagonal identically; no clipping or post-hoc renormalisation is involved. For a single channel with $\Omega=1$ and $v=e_i+e_j$, writing $a:=\sigma_W$, $p:=C^{-}_{ij}$, $X:=C^{-}_{ii}$ and $Y:=C^{-}_{jj}$, the block entry moves by
\[
\Delta_{R,ij}=\frac{p+a}{\sqrt{(X+a)(Y+a)}}-\frac{p}{\sqrt{XY}},
\]
whose derivative at $a=0$ is $(XY)^{-1/2}\,[1-p(X+Y)/(2XY)]$. The block entry therefore rises near the wall precisely when $p<2XY/(X+Y)$, the harmonic mean of the two variances. By the arithmetic--geometric mean inequality that mean is at most $\sqrt{XY}$, with equality exactly when $X=Y$, and $|p|\le\sqrt{XY}$ because $C$ is positive semidefinite; so the condition holds for every admissible $p$ when the variances are equal, and can fail only for a pair that is both strongly correlated and strongly heteroscedastic. The leading coefficient is not $a/\sqrt{XY}$ in general.

A pair with one coin inside the block and one outside behaves differently in sign: because only the inside coin's diagonal moves, $\Delta_{R,ik}=-aC^{-}_{ik}/(2(X)^{3/2}\sqrt{C^{-}_{kk}})+O(a^2)$, which is opposite-signed to $C^{-}_{ik}$ and vanishes exactly when $C^{-}_{ik}=0$. Its size relative to the block entry is of order $|C^{-}_{ik}|/(2C^{-}_{ii})$, small when the pre-event cross-covariance is small relative to the block coin's variance but not bounded away from the block term in general.

Under the midpoint calibration $\Sigma_G(\theta_W)=\tfrac12(\Sigma^{+}+\Sigma^{-})$ at each wall, a Gaussian copula surface $\Sigma_G$ satisfies $\Sigma(\theta^{\epsilon}_{\pm})-\Sigma_G(\theta^{\epsilon}_{\pm})=\pm\tfrac12\Delta_R+O(\epsilon)$ as the wall is approached. The magnitude is a consequence of the midpoint calibration rather than of the Gaussian family; what extends to any smooth surface needs only local continuity - a smooth surface cannot track a function that is constant on chambers and steps between them, so its distance from $\Sigma$ is of the order of the jump and does not vanish.

\subsection*{Panel and windows}
The panel is 28 stablecoins with daily coverage dense enough to enter the panel, drawn from the DeFiLlama public archive over 2020--2025 and cross-checked against Chainlink oracle feeds and CoinGecko quotations; no proprietary or issuer-internal data is used. The pre-event window is the 60 calendar days preceding the event, trimmed to observations with non-missing prices (typically 46 trading days), and ends before the acute window begins; the acute window is the eleven calendar days centred on the event. Because the acute window is centred it contains days on both sides of the crossing, which pulls $\rho^{+}$ toward $\rho^{-}$ and biases $\Delta\rho$ toward zero: the test is conservative in this respect. No pre-window observation enters the definition of $B_W$.

For each event and each pair with at least five usable observations in each window, $\Delta\rho_{ij}:=\rho^{+}_{ij}-\rho^{-}_{ij}$, with $\rho$ the Pearson correlation of daily peg deviations. This is the empirical counterpart of $\Delta_{R,ij}$: the pre-window estimates the pre-crossing chamber value, the acute window the post-crossing value. The identification of a pre-event window with a chamber requires the window to be long enough for the peg-deviation law to settle inside the regime and short enough to precede the event; the first is examinable through the window grid, the second we take as given.

\subsection*{Test and inference}
Pairs are classified BLOCK (both coins in $B_W$), CROSS (exactly one) and OUTSIDE (neither). The tested criterion is that the BLOCK layer mean exceeds both others. The comparison between CROSS and OUTSIDE is reported as an empirical quantity and is not used as a test: the pure crossing law places CROSS below OUTSIDE, since the cross-pair image is opposite-signed, but the observed level of every layer is contaminated by crisis-wide co-movement in proportions this design cannot separate.

A crisis raises co-movement across the whole panel, so the substantive statistic is the block-versus-rest contrast $\widehat{\mathrm{DiD}}=\overline{\Delta\rho}_{\mathrm{BLOCK}}-\overline{\Delta\rho}_{\mathrm{BLOCK}^{c}}$. This is valid only under an identifying condition, which we state explicitly: Assumption (H3) implies that pairs outside the block carry no crossing contribution, so the theory predicts a zero OUTSIDE layer, whereas the observed OUTSIDE means are not zero. A component the theory does not model is therefore present, and $\widehat{\mathrm{DiD}}$ identifies the block effect only if that component is additive and does not differ systematically across layers. We adopt that as a condition, not as a consequence.

Primary inference is a coin-level permutation: the block label is redrawn over sets of coins of the same size as $B_W$ and $\widehat{\mathrm{DiD}}$ recomputed, with 20{,}000 draws, which asks how unusual the true block is among the $\binom{n}{|B_W|}$ coin sets that could have been so labelled. Pair-level permutation, which treats pairs as exchangeable, is reported alongside as the anti-conservative comparison. A coin-level bootstrap interval is also reported; it is conditional on the block assignment, so it measures sampling variability of coins under that assignment rather than the distinctiveness of the block.

\subsection*{Window grid and robustness}
The canonical specification is a 60-day pre-window with a $\pm5$-day acute window, fixed by the panel-construction script before the test was designed. A $3\times3$ grid crosses pre-windows of 30, 60 and 90 days with acute half-widths of 3, 5 and 10 days. The script that produces the grid re-derives the canonical row exactly, which is the check that the grid is comparable across specifications.

\subsection*{The amplitude}
The scale $\sigma_W$ is left free. Because the normalising map is a nonlinear congruence it destroys the rank-one form of the jump, so what survives into correlation space is sign structure and support rather than scale; a test built on those is invariant to $\sigma_W$ at leading order in $a$, which is the order at which the criteria are stated. Fixing $\sigma_W$ would matter only for a test of magnitude, which requires more multi-coin-block events than the record contains.

\subsection*{Data and code availability}
The panel is assembled from public archives (DeFiLlama, FRED, Yahoo Finance, Etherscan) with source URLs, retrieval timestamps, sizes and SHA-256 checksums recorded for every raw file. The reproducibility bundle contains the LaTeX source, all figure-generation and audit scripts, the raw data archive and a checksummed manifest, and the pipeline completes in under three minutes. Code is available at https://github.com/dyzchina/crypto4.
